\clearpage

\addcontentsline{toc}{chapter}{Abstract}
\begin{center}
	\textbf{\large Abstract}
\end{center}

Named Entity Recognition (NER) in specialized domains requires costly manual annotation and a predefined type inventory. This thesis addresses \textbf{open-world NER}: detecting entity mentions from unlabeled text and automatically discovering and naming latent entity types without any target-domain supervision.

Building on the OWNER framework, we propose \textbf{Span2Type} with three targeted improvements. First, we replace the BIO mention detector with a multi-view divide-and-transfer (DTrans) encoder~\cite{zhang-etal-2024-dtrans} that jointly trains BIO, Start-End, and Tie-Break decoding heads, improving cross-domain span recall. Second, we introduce soft-prompt entity encoding: trainable virtual tokens are prepended to a fill-in-the-blank BERT template, and the \texttt{[MASK]}-position embedding serves as a type-discriminative representation for clustering. Third, we propose MMR-based exemplar selection for cluster naming, querying either an MLM or a local LLM with diverse cluster evidence to generate concise, human-readable type labels.

Evaluated on five CrossNER domains, Span2Type achieves an average AMI of 61.4, outperforming the OWNER baseline by 12.0 points, and produces cluster names with BERTScore-F1 of 0.924.

\vspace*{1cm}

\noindent\textbf{Keywords:} Open-world Named Entity Recognition; Unsupervised NER; Multi-view Span Detection; Soft Prompt Entity Typing; Cluster Naming

\begin{comment}


\end{comment}


